\setchapterpreamble[u]

\chapter{Introduction} \label{chap:Introduction}
\minitoc\vspace{1em}
\blindtext

\todo{add some TODOs which are visible as marginal notes in the PDF}{and some comment for the todo} This is how to reference \figref{fig:SWC}. \gls{software testing} is the art of \gls{testing}.

\blindtext

\begin{figure}[ht]
  \centering
    \includegraphics[width=0.7\textwidth]{rwth_swc_cmyk.pdf}
  \caption{SWC} \label{fig:SWC}
\end{figure}

\Blindtext See \tabref{tab:Price list}. \comment{foo bar}{... and here a comment}

\begin{table}
  \begin{center}
    \begin{tabular}{llr}
      \toprule
      \multicolumn{2}{c}{Item} \\
      \cmidrule(r){1-2}
      Animal & Description & Price (\$) \\
      \midrule
      Gnat      & per gram & 13.65 \\
                & each     &  0.01 \\
      Gnu       & stuffed  & 92.50 \\
      Emu       & stuffed  & 33.33 \\
      Armadillo & frozen   & 8.99  \\
      \bottomrule
    \end{tabular}
  \end{center}
  \caption{Price list} \label{tab:Price list}
\end{table}

\blindtext

Defining \lstinlineJava{class MyClass} in \lstref{lst:MyClass.java}:

%\begin{lstlisting}[caption={MyClass.java}, label={lst:MyClass.java}, language=java, style=eclipse]
\begin{lstlistingJava}{MyClass.java}
public class MyClass {
    int    myVar1     = 27;
    float  thisIsCool = 3.234f;
    String myString   = "Hello World!";

    /** This is a Javadoc Comment */
    /* This is a normal comment */
    public MyClass(String param) {
        myString = param; // comment
    }
}
\end{lstlistingJava}

The style of highlighting is the same as used by Eclipse \cite{Eclipse} \ldots

In \chapref{chap:Background}, we will give the background for this thesis... In \secpref{sec:Environment}, foo bar \gls{testing}\index{testing} \cite{ElectronicsEngingeers1990}. Using \gls{design pattern}s in \gls{software architecture}.
